% \iffalse meta-comment
%
% hit-subfigure.dtx 是各个类共用的 subfigure 兼容层
%
% \fi
%
% \section{共享 subfigure 兼容层}
%
% \changes{v3.2a}{2026/08/07}{\pkg{subfigure} 兼容层抽出来给各个类共用，
%   原先 \file{hithesisbook.cls} 与 \file{hithesisart.cls} 各存一份，290 行逐字节相同}
% \pkg{subfigure} 二十多年没人维护，与 \pkg{cleveref}、\pkg{hyperref} 配合有问题，
% v3.2a 换成了 \pkg{subcaption}。这个文件把旧接口在新实现上重做一遍，让写着
% \cs{subfigure}、\cs{subtable}、\cs{ifsubcaphang} 的老文档继续能编。
%
% \textbf{整个文件是一次性的。} 里面没有一行是模板自己要用的东西——\cs{subfigure} 的
% 六个长度、\cs{bisubcaption}、五个开关的转译、\cs{subfigure} 与 \cs{subtable} 的
% 重定义，全都只为接住旧写法。兼容期在 v4（预计~2027~年发布）结束，那时把
% \file{hithesis.ins} 里这三行删掉、连同本文件一起移除即可，不必去别的文件里挑拣。
% 移除前记得把示例正文改写成 \cs{subcaptionbox}，那里现在用的还是旧写法。
%
% 子图题的字号与标签分隔（\cs{subcapsize}、\cs{DeclareCaptionFont}、
% \cs{captionsetup}\oarg{sub}）不在这里：那是版式，各个类自己定，与旧接口无关。
%
%    \begin{macrocode}
%<*shared-subfigure>
%    \end{macrocode}
%
% 绝大部分是普通定义，直接放在文件顶层——包进一个函数的话，里面十几个带参数的
% 定义都要把 \texttt{\#} 加倍，读起来与出错风险都不划算。
%
% 六个长度只在这里声明，赋值在 \cs{hit@setup@subfigure} 里。
%    \begin{macrocode}
\skip_new:N \subfigtopskip
\skip_new:N \subfigcapskip
\skip_new:N \subfigbottomskip
\dim_new:N  \subfigcaptopadj
\dim_new:N  \subfigcapmargin
\skip_new:N \subfiglabelskip
%    \end{macrocode}
%
% \cs{subcaptionbox} 要求 \cs{label} 写在图题里，而旧写法把它写在内容里
% （\pkg{subfigure} 自己 \cs{let} 掉 \cs{label} 来接）。这里照样接住：内容先
% 装箱，装箱时把 \cs{label} 的参数记下来，再塞回图题参数末尾。
%    \begin{macrocode}
\box_new:N \l__hit_subfloat_content_box
\box_new:N \l__hit_subfloat_caption_box
\box_new:N \l__hit_subfloat_measure_box
\dim_new:N \l__hit_subfloat_total_dim
\tl_new:N \g__hit_subfloat_label_tl
\cs_new_eq:NN \__hit_subfloat_label:n \label
\cs_new_protected:Npn \__hit_subfloat:nn #1#2
  {
    \tl_gclear:N \g__hit_subfloat_label_tl
    \hbox_set:Nn \l__hit_subfloat_content_box
      {
        \cs_set_protected:Npn \label ##1 { \tl_gset:Nn \g__hit_subfloat_label_tl {##1} }
        #2
      }
%    \end{macrocode}
%
% 图题先单独装进 \cs{vbox} 再放进竖直列表。\pkg{caption} 排图题时会先放一条零高
% \cs{hrule}（\cs{vspace} 在竖直模式下的常规动作），\cs{hrule} 会掐掉 \hologo{TeX}\ 的
% 行间胶，图与图题之间就少了一整个 \cs{baselineskip}。装箱之后 \cs{vbox} 的高度
% 就是首行基线的位置，行间胶照常算，与 \pkg{subfigure} 的结果一致。
%    \begin{macrocode}
    \tl_if_empty:NF \g__hit_subfloat_pending_label_tl
      {
        \tl_gset:Nx \g__hit_subfloat_label_tl { \g__hit_subfloat_pending_label_tl }
        \tl_gclear:N \g__hit_subfloat_pending_label_tl
      }
    \dim_compare:nNnT { \box_wd:N \l__hit_subfloat_content_box } = { 0pt }
      { \box_set_wd:Nn \l__hit_subfloat_content_box { \linewidth } }
%    \end{macrocode}
%
% \pkg{subfigure} 单行子图题排成裸 \cs{hbox}，高深随文字；排不下才退回
% \cs{parbox}，那时才有 strut。\pkg{caption} 一律带 strut，深度固定，图之后的
% 材料会高一截。照 \pkg{subfigure} 的做法先拿同样的字体量一遍宽度，装得下就
% 关掉 strut。量的时候计数器还没走，\cs{thesubfigure} 拿到的是上一个号，宽度
% 一样，不影响判断。
%    \begin{macrocode}
    \hbox_set:Nn \l__hit_subfloat_measure_box
      {
        \subcapsize \thesubfigure \hskip \subfiglabelskip #1
      }
    \tex_setbox:D \l__hit_subfloat_caption_box \tex_vbox:D
      {
        \hsize \box_wd:N \l__hit_subfloat_content_box
        \linewidth \hsize
        \captionsetup [ sub ] { strut = false }
        \bool_if:NT \g__hit_subfloat_reuse_bool
          {
            \addtocounter { sub \@captype } { -1 }
            \bool_gset_false:N \g__hit_subfloat_reuse_bool
          }
        \bool_gset_false:N \g__hit_subfloat_skip_label_bool
        \subcaption
          {
            #1
            \bool_if:NF \g__hit_subfloat_skip_label_bool
              {
                \tl_if_empty:NF \g__hit_subfloat_label_tl
                  { \exp_args:NV \__hit_subfloat_label:n \g__hit_subfloat_label_tl }
              }
          }
      }
%    \end{macrocode}
%
% \pkg{subfigure} 的子表题在表上方（\cs{ifsubtabletopcap}），子图题在图下方。
% 盒子是这里手工摆的，\pkg{caption} 的 \option{position} 管不到，按 \cs{@captype}
% 自己分。上方时上下留白对调，并按 \pkg{subfigure} 的做法补 \cs{subfigcaptopadj}。
%
% \pkg{subfigure} 多行子图题用 \cs{parbox}\texttt{[t]}：高度是首行的高度，其余算深度。
% \cs{vbox} 相反，把一切折进高度、深度为 0，放进竖直列表时参考点就错了。这里拿
% 量宽度那个 \cs{hbox} 的高度当首行高度，把图题盒子的高深重设一遍。单行时结果与
% \pkg{subfigure} 的裸 \cs{hbox} 逐值相同。
%    \begin{macrocode}
    \tl_gclear:N \g__hit_subfloat_label_tl
    \dim_set:Nn \l__hit_subfloat_total_dim
      { \box_ht:N \l__hit_subfloat_caption_box + \box_dp:N \l__hit_subfloat_caption_box }
    \box_set_ht:Nn \l__hit_subfloat_caption_box { \box_ht:N \l__hit_subfloat_measure_box }
    \box_set_dp:Nn \l__hit_subfloat_caption_box
      { \l__hit_subfloat_total_dim - \box_ht:N \l__hit_subfloat_measure_box }
    \leavevmode
    \str_if_eq:eeTF { \@captype } { table }
      {
        \vtop
          {
            \hsize \box_wd:N \l__hit_subfloat_content_box
            \vbox
              {
                \vskip \subfigbottomskip
                \box_use_drop:N \l__hit_subfloat_caption_box
              }
            \vskip \subfigcapskip
            \vskip \subfigcaptopadj
            \box_use_drop:N \l__hit_subfloat_content_box
            \vskip \subfigtopskip
          }
      }
      {
        \vtop
          {
            \hsize \box_wd:N \l__hit_subfloat_content_box
            \vbox { \vskip \subfigtopskip \box_use_drop:N \l__hit_subfloat_content_box }
            \vskip \subfigcapskip
            \box_use_drop:N \l__hit_subfloat_caption_box
            \vskip \subfigbottomskip
          }
      }
  }
%    \end{macrocode}
%
% \pkg{subcaption} 把 \texttt{subfigure} 和 \texttt{subtable} 定义成环境，而
% \cs{begin}\marg{subfigure} 展开出来就是 \cs{subfigure}，跟旧的命令写法同名。
% 两套语法要并存，只能在入口按 \cs{@currenvir} 分流：从 \cs{begin} 进来的走
% \pkg{subcaption} 的环境，直接调用的走兼容壳。
%
% \begin{macro}{\hit@setup@subfigure}
% 两类东西要等到调用时才做。一是接管 \cs{subfigure} 与 \cs{subtable}，那得先等
% \pkg{subcaption} 加载完、把原定义存下来。二是六个长度的赋值：\cs{subfiglabelskip}
% 的值是 \texttt{0.33em}，跟着字体走，在共享层加载时字体还没配好，\texttt{em} 换算
% 基准不对——踩过一次，子图标签后面的空格没了，\texttt{(a) 打高尔夫球} 变成
% \texttt{(a)打高尔夫球}。声明留在顶层，赋值放这里，时机与合并前一致。
%    \begin{macrocode}
\cs_new_protected:Npn \hit@setup@subfigure
  {
    \skip_set:Nn \subfigtopskip    { 5pt }
    \skip_set:Nn \subfigcapskip    { 0pt }
    \skip_set:Nn \subfigbottomskip { 5pt }
    \dim_set:Nn  \subfigcaptopadj  { 3pt }
    \dim_set:Nn  \subfigcapmargin  { 0pt }
    \skip_set:Nn \subfiglabelskip  { 0.33em~plus~0.07em~minus~0.03em }
    \cs_if_exist:NF \subcapsize { \cs_new:Npn \subcapsize { \wuhao } }
    \cs_new_eq:NN \__hit_subfigure_environment: \subfigure
    \cs_new_eq:NN \__hit_subtable_environment: \subtable
    \cs_set_protected:Npn \subfigure
      {
        \str_if_eq:eeTF { \@currenvir } { subfigure }
          { \__hit_subfigure_environment: }
          { \__hit_subfloat_command: }
      }
    \cs_set_protected:Npn \subtable
      {
        \str_if_eq:eeTF { \@currenvir } { subtable }
          { \__hit_subtable_environment: }
          { \__hit_subfloat_command: }
      }
  }
%    \end{macrocode}
%
% 旧写法有三个习惯用法，兼容期内都要接住：
% \begin{itemize}
% \item \cs{subfigure}\marg{\cs{label}\marg{键}}：内容是空的，只为了把标签挂到
%   下一个编号上。这里记下标签留给下一个子图，本身什么也不排。
% \item 紧跟其后的 \cs{addtocounter}\marg{subfigure}\marg{-2}：那是给上面那招凑
%   计数器用的，\pkg{subcaption} 的计数器语义不同，直接吃掉。
% \item \cs{subfigure}\oarg{英文}\marg{\cs{subfigure}\oarg{中文}\marg{图}}：双语
%   子图题的套娃，拆开转成 \cs{bisubcaption}。
% \end{itemize}
%    \begin{macrocode}
\tl_new:N \g__hit_subfloat_pending_label_tl
\cs_new_protected:Npn \__hit_subfloat_discard:w \addtocounter #1#2
  { \str_if_eq:nnF {#1} { subfigure } { \addtocounter {#1} {#2} } }
\cs_new_protected:Npn \__hit_subfloat_discard_rewind:
  {
    \peek_remove_spaces:n
      { \peek_meaning:NT \addtocounter { \__hit_subfloat_discard:w } }
  }
%    \end{macrocode}
%
% 双语子图题照旧写法的结构递归摆一次：内层是「图 + 中文题」，外层把它整个当内容
% 再排一次英文题。两行之间那段 \cs{subfigbottomskip} 是这么来的，在图题文本里补
% 不出来。外层排图题前把计数器退一格，两行才共用同一个编号，与旧写法一致。
%    \begin{macrocode}
\bool_new:N \g__hit_subfloat_reuse_bool
\bool_new:N \g__hit_subfloat_skip_label_bool
\cs_new_protected:Npn \__hit_subfloat_bicaption:nnn #1#2#3
  {
    \__hit_subfloat:nn {#2}
      {
        \bool_gset_true:N \g__hit_subfloat_skip_label_bool
        \__hit_subfloat:nn {#1} {#3}
        \bool_gset_true:N \g__hit_subfloat_reuse_bool
      }
  }
%    \end{macrocode}
%
% \changes{v3.2a}{2026/08/07}{旧写法加废弃提示（\issue[284]）}
% 旧写法的提示。每种只报一次，用 \cs{\_\_hit\_warn\_once:nnn} 记名。\pkg{subfigure}
% 专有的开关与命令没有对应写法，报出来总比「未定义控制序列」强。
%    \begin{macrocode}
\cs_new_protected:Npn \__hit_subfloat_dropped:nn #1#2
  { \__hit_warn_once:nnn {#1} { removed-subfig } {#2} }
\cs_new_protected:Npn \__hit_subfloat_switch:nnnn #1#2#3#4
  {
    \cs_new_protected:cpn { #1 true }
      {
        \__hit_warn_once:nnn { if#1 } { mapped-subfig } { #2 = #3 }
        \captionsetup [ sub ] { #2 = #3 }
      }
    \cs_new_protected:cpn { #1 false }
      {
        \__hit_warn_once:nnn { if#1 } { mapped-subfig } { #2 = #4 }
        \captionsetup [ sub ] { #2 = #4 }
      }
  }
\__hit_subfloat_switch:nnnn { subcaphang }
  { format } { hang } { plain }
\__hit_subfloat_switch:nnnn { subcapcenter }
  { justification } { centering } { justified }
\__hit_subfloat_switch:nnnn { subcapcenterlast }
  { justification } { centerlast } { justified }
\__hit_subfloat_switch:nnnn { subcapnooneline }
  { singlelinecheck } { false } { true }
\__hit_subfloat_switch:nnnn { subcapraggedright }
  { justification } { raggedright } { justified }
\cs_new_protected:Npn \listsubcaptions
  { \__hit_subfloat_dropped:nn { listsubcaptions } { captionsetup[sub]{list=true} } }
%    \end{macrocode}
%
% 拆嵌套写法之前先提示。
%    \begin{macrocode}
\cs_new_protected:Npn \__hit_subfloat_unnest:nw #1 \subfigure [#2]#3 \q_stop
  {
    \__hit_warn_once:nnn { nested-subfigure } { deprecated-subfig }
      { 嵌套的~subfigure~写法已不推荐。 }
    \__hit_subfloat_bicaption:nnn {#2} {#1} {#3}
  }
\cs_new_protected:Npn \__hit_subfloat_bicaption:w \bisubcaption #1#2 \q_stop #3
  { \__hit_subfloat_bicaption:nnn {#1} {#2} {#3} }
\cs_new_protected:Npn \__hit_subfloat_no_caption:n #1
  {
    \tl_gclear:N \g__hit_subfloat_label_tl
    \hbox_set:Nn \l__hit_subfloat_content_box
      {
        \cs_set_protected:Npn \label ##1 { \tl_gset:Nn \g__hit_subfloat_label_tl {##1} }
        #1
      }
    \dim_compare:nNnTF { \box_wd:N \l__hit_subfloat_content_box } = { 0pt }
      {
        \tl_gset_eq:NN \g__hit_subfloat_pending_label_tl \g__hit_subfloat_label_tl
        \__hit_warn_once:nnn { empty-subfigure } { deprecated-subfig }
          { 用空的~subfigure~占标签已不推荐。 }
        \leavevmode
        \__hit_subfloat_discard_rewind:
      }
      { \__hit_subfloat:nn { } { \box_use_drop:N \l__hit_subfloat_content_box } }
  }
\NewDocumentCommand \__hit_subfloat_command: { o m }
  {
    \IfNoValueTF {#1}
      { \__hit_subfloat_no_caption:n {#2} }
      {
        \tl_if_head_eq_meaning:nNTF {#2} \subfigure
          { \__hit_subfloat_unnest:nw {#1} #2 \q_stop }
          {
            \tl_if_head_eq_meaning:nNTF {#1} \bisubcaption
              { \__hit_subfloat_bicaption:w #1 \q_stop {#2} }
              {
                \tl_if_in:nnT {#2} { \subfigure }
                  {
                    \__hit_warn_once:nnn { unrecognised-nesting }
                      { deprecated-subfig }
                      { 内容里有~subfigure，但不在最前面，认不出是双语写法。 }
                  }
                \__hit_subfloat:nn {#1} {#2}
              }
          }
      }
  }
%    \end{macrocode}
%
% \changes{v3.2a}{2026/08/07}{加 \cs{bisubcaption}，双语子图题不再靠嵌套
%   \cs{subfigure} 凑（\issue[284]）}
% 双语子图题。它只是个记号，由 \cs{subfigure} 认出来自己拆：两行图题之间隔着一段
% \cs{subfigbottomskip}，那是旧写法里内外两层子浮动体叠出来的结构，在图题文本里
% 补不进去（\cs{vskip} 与 \cs{\\} 在量宽度的 \cs{hbox} 里都非法）。所以拆开之后
% 照旧写法的结构递归摆一遍。
%    \begin{macrocode}
\cs_new_protected:Npn \bisubcaption #1#2 { }
%    \end{macrocode}
% \end{macro}
%
%    \begin{macrocode}
%</shared-subfigure>
%    \end{macrocode}
%
%
% \Finale
